\section{Flat base change for the pushforward of a quasi-coherent sheaf ($i=0$)}
\label{mk-chap:Cohomology_FlatBaseChange}

\section{Two routes to base change, and which one is active}
\label{mk-sec:fbc_route_decision}

There are two distinct ways to express ``the pushforward commutes with flat base
change'' at the level of \(\mathcal{O}_{S'}\)-modules. This chapter records both, but
commits to one as the load-bearing construction.

\emph{The canonical-mate route (Stacks-faithful, superseded).} One fixes
the canonical base-change morphism
\(\beta : g^*(f_*\mathcal{F}) \to f'_*((g')^*\mathcal{F})\) of
Definition~\ref{mk-def:pushforward_base_change_map} --- the adjoint mate of the
\(((g')^*,(g')_*)\)-unit --- and proves \(\mathrm{IsIso}\,\beta\) directly,
affine-locally. This is the route the Stacks Project takes, and it remains the
mathematically canonical formulation, but its closure requires a \(2\)-categorical
mate-unwinding (identifying the abstract mate, transported through the four affine
dictionaries, with the concrete cancellation isomorphism
\(\operatorname{cancelBaseChange}\)) that has no known direct proof. Only the
definition of \(\beta\) is retained; the route's obligation lemmas and its two
theorem statements were removed as superseded once the concrete-tilde route below
closed the downstream consumers.

\emph{The concrete-tilde construction (active, consumed downstream).} Rather than prove
that the \emph{canonical} map is an isomorphism, one \emph{constructs a fresh} natural
isomorphism
\[
  e : g^* \circ f_* \;\cong\; f'_* \circ (g')^*
\]
directly, by gluing the affine-local isomorphisms of
Lemma~\ref{mk-lem:affine_pushforward_pullback_baseChange} (each assembled from the two
tilde dictionaries and Mathlib's \(\operatorname{cancelBaseChange}\)) along the affine
basis of \(S'\). This construction \textbf{never forms the adjoint mate}, so it
entirely sidesteps the mate\(\,\leftrightarrow\,\operatorname{cancelBaseChange}\)
coherence obligation of the superseded canonical route: the only remaining
content is the naturality of the cancellation chain across affine restrictions of
\(S'\) (Lemma~\ref{mk-lem:pushforwardPullbackBaseChange_restrict_naturality}), a concrete
localization-and-tensor compatibility rather than a mate identification.

\emph{Why the active route suffices.} Everything downstream of this chapter --- the
{\v C}ech tensorial base-change isomorphism (the cosimplicial residual \(e\) of the
{\v C}ech chapter) and, through it, the Kleiman-style representability inputs ---
consumes \emph{only} a base-change isomorphism (a \(\mathrm{Nonempty}\) / \(\cong\)
datum) that commutes with the restriction maps, \emph{not} the assertion that the
canonical map \(\beta\) is that isomorphism. The fresh \(e\) supplies exactly such an
isomorphism, and by construction it commutes with restriction (it is glued from
restriction-compatible affine pieces). The canonical map
\(\operatorname{pushforwardBaseChangeMap}\) is not a protected declaration, so
re-plumbing the downstream consumers to receive \(e\) is permitted. Consequently the
key constituent of this chapter is
Lemma~\ref{mk-lem:affine_pushforward_pullback_baseChange} and the gluing chain
(\S\ref{mk-sec:fbc_gluing}) built on it; the canonical-mate nodes are retained for
faithfulness but are off the critical path.

\section{The canonical base-change map}

\begin{definition}[Base-change map for the pushforward]\label{mk-def:pushforward_base_change_map}
  \dcref{custom}

  \textit{Source: Stacks Project, Cohomology of Schemes, \S\,Cohomology and base
  change, I.}
  Let \(f : X \to S\) be a morphism of schemes, \(\mathcal{F}\) a quasi-coherent
  \(\mathcal{O}_X\)-module, and \(g : S' \to S\) an arbitrary morphism, giving the
  cartesian square above with projection \(g' : X' \to X\) and
  \(\mathcal{F}' = (g')^*\mathcal{F}\). The \emph{base-change map for the
  pushforward} is the canonical morphism of \(\mathcal{O}_{S'}\)-modules
  \[
    g^*\bigl(f_*\mathcal{F}\bigr) \longrightarrow f'_*\bigl((g')^*\mathcal{F}\bigr)
    = f'_*\mathcal{F}' .
  \]
  It is adjoint to the morphism
  \(f_*\mathcal{F} \to g_*\, f'_*\mathcal{F}' = f_*\,(g')_*\mathcal{F}'\) obtained
  by applying \(f_*\) to the unit
  \(\mathcal{F} \to (g')_*(g')^*\mathcal{F}\) of the
  \(((g')^*, (g')_*)\)-adjunction, using the commutativity
  \(g \circ f' = f \circ g'\) of the square.
\end{definition}

\section{The affine case}

\subsection{Locality of isomorphisms for \texttt{Scheme.Modules} morphisms}

Before treating the affine computation we record three project-local supplements
to Mathlib. Mathlib provides the per-open criterion
(\(\mathrm{IsIso}\,\varphi \iff \forall U,\ \mathrm{IsIso}(\varphi_U)\)) and a
stalkwise isomorphism criterion for \(\mathrm{TopCat.Sheaf}\)-valued morphisms, but
it does not package the stalk-local or basis-local criterion at the level of
morphisms of \(\mathcal{O}_X\)-modules. The following lemmas bridge that gap; they
are exactly the locality tools used in the first reduction of the affine
base-change lemma below, where one checks the base-change map after restricting to
the affine opens of the base. As bespoke infrastructure they carry no external
citation and stand on their proof sketches.

\begin{lemma}[Stalk-local criterion for isomorphisms of \(\mathcal{O}_X\)-modules]\label{mk-lem:modules_isIso_iff_stalk}
  \dcref{custom}

  Let \(X\) be a scheme and \(\varphi : M \to N\) a morphism of sheaves of
  \(\mathcal{O}_X\)-modules. Then \(\varphi\) is an isomorphism if and only if the
  underlying morphism of abelian-group presheaves induces an isomorphism on the
  stalk at every point \(x \in X\).
\end{lemma}

\begin{proof}


If \(\varphi\) is an isomorphism, then so is its image under the forgetful
  functor to abelian-group presheaves, and the stalk functor at each \(x\),
  being a functor, carries isomorphisms to isomorphisms. Conversely, assume the
  stalk map is an isomorphism at every \(x\). Package the underlying abelian-group
  presheaves of \(M\) and \(N\) as sheaves of abelian groups, and \(\varphi\) as a
  morphism of these sheaves; the stalkwise-isomorphism hypothesis lets us apply the
  \(\mathrm{TopCat.Sheaf}\)-level stalk-local isomorphism criterion to conclude that
  this morphism of sheaves of abelian groups is an isomorphism. Pushing it through
  the forgetful functor to presheaves yields an isomorphism of the underlying
  abelian-group presheaf morphisms, and since the forgetful functor from
  \(\mathcal{O}_X\)-modules to abelian-group presheaves reflects isomorphisms,
  \(\varphi\) itself is an isomorphism.
\end{proof}

\begin{lemma}[Basis-local criterion for isomorphisms of \(\mathcal{O}_X\)-modules]\label{mk-lem:modules_isIso_of_isBasis}
  \dcref{custom}

  \uses{mk-lem:modules_isIso_iff_stalk}
  Let \(X\) be a scheme, let \(\{B_i\}_{i \in \iota}\) be a family of opens whose
  range is a basis of the topology of \(X\), and let \(\varphi : M \to N\) be a
  morphism of sheaves of \(\mathcal{O}_X\)-modules. If \(\varphi\) restricts to an
  isomorphism on the sections over every basic open \(B_i\), then \(\varphi\) is an
  isomorphism.
\end{lemma}

\begin{proof}
  \uses{mk-lem:modules_isIso_iff_stalk}
  By Lemma~\ref{mk-lem:modules_isIso_iff_stalk} it suffices to prove that the induced
  stalk map is bijective at each point \(x \in X\). For injectivity, a germ that
  vanishes at \(x\) already vanishes on some open neighbourhood, which may be
  shrunk to a basic open \(B_i\); on \(B_i\) the section map \(\varphi_{B_i}\) is
  injective, so the original section is zero, giving injectivity of the stalk map
  from injectivity over the basis. For surjectivity, an arbitrary germ at \(x\) is
  represented by a section of \(N\) over some basic open \(B_i\) containing \(x\);
  since \(\varphi_{B_i}\) is surjective, this section lifts to a section of \(M\)
  over \(B_i\), whose germ at \(x\) maps to the given germ. Thus the stalk map is
  bijective, and \(\varphi\) is an isomorphism.
\end{proof}

\begin{lemma}[Affine-open locality criterion for isomorphisms of \(\mathcal{O}_X\)-modules]\label{mk-lem:modules_isIso_iff_affineOpens}
  \dcref{custom}

  \uses{mk-lem:modules_isIso_of_isBasis}
  Let \(X\) be a scheme and \(\varphi : M \to N\) a morphism of sheaves of
  \(\mathcal{O}_X\)-modules. Then \(\varphi\) is an isomorphism if and only if it
  restricts to an isomorphism on the sections over every affine open of \(X\).
\end{lemma}

\begin{proof}
  \uses{mk-lem:modules_isIso_of_isBasis}
  The forward implication is immediate, since the section functor over a fixed open
  preserves isomorphisms. For the converse, the affine opens of a scheme form a
  basis of its topology, so the claim is the special case of
  Lemma~\ref{mk-lem:modules_isIso_of_isBasis} in which the basis is the family of
  affine opens.
\end{proof}

\subsection{The affine computation}

The whole result rests on the following elementary computation in the affine
setting, where the base-change map is visibly an isomorphism because tensoring is
associative. Before stating it we isolate the one ingredient not yet available as
a ready-made result on which the affine reduction turns: the description of the
pushforward of a tilde-module along a morphism of affine schemes.

We first record three bespoke \(\Gamma\)-fragment supports that compute the global
sections of an affine pushforward; they carry no external citation and stand on
their proof sketches.

\begin{lemma}[Naturality of the global-sections comparison square]\label{mk-lem:globalSectionsIso_hom_comp_specMap_appTop}
  \dcref{custom}

  Let \(\varphi : R \to R'\) be a homomorphism of commutative rings, and write
  \(\mathrm{gs}_R : \Gamma(\operatorname{Spec} R, \top) \cong R\) and
  \(\mathrm{gs}_{R'} : \Gamma(\operatorname{Spec} R', \top) \cong R'\) for the
  global-sections isomorphisms of the structure sheaves. Then the ring square
  \[
    \mathrm{gs}_R \;\circ\; (\operatorname{Spec}\varphi)^{\sharp}_{\top}
    \;=\; \varphi \;\circ\; \mathrm{gs}_{R'}
  \]
  commutes, where \((\operatorname{Spec}\varphi)^{\sharp}_{\top}\) is the
  top-open component of the structure-sheaf comparison map underlying
  \(\operatorname{Spec}\varphi\).
\end{lemma}

\begin{proof}


This is the naturality of the unit/counit of the
  \(\Gamma \dashv \operatorname{Spec}\) adjunction read off at the top open: the
  global-sections map of \(\operatorname{Spec}\varphi\) is, under the
  global-sections isomorphisms on each side, conjugate to \(\varphi\) itself.
  Transposing the naturality square of the \(\operatorname{Spec}\) unit gives
  exactly the asserted ring equation.
\end{proof}

\begin{lemma}[Global sections of an affine pushforward]\label{mk-lem:gammaPushforwardIso}
  \dcref{custom}

  \uses{mk-lem:globalSectionsIso_hom_comp_specMap_appTop}
  Let \(\varphi : R \to R'\) be a homomorphism of commutative rings and let
  \(N\) be an \(\operatorname{Spec}(R')\)-module. Then there is a canonical
  isomorphism of \(R\)-modules
  \[
    \Gamma\bigl((\operatorname{Spec}\varphi)_* N\bigr)
    \;\cong\;
    \operatorname{restr}_\varphi\bigl(\Gamma N\bigr),
  \]
  identifying the \(R\)-module of global sections of the pushforward with the
  restriction of scalars along \(\varphi\) of the \(R'\)-module of global sections
  of \(N\).
\end{lemma}

\begin{proof}
  \uses{mk-lem:globalSectionsIso_hom_comp_specMap_appTop}
  Because \((\operatorname{Spec}\varphi)^{-1}(\top) = \top\), the global sections of
  the pushforward \((\operatorname{Spec}\varphi)_* N\) are, on underlying abelian
  groups, the global sections of \(N\) with no transport along the open-set
  identification. The two sides therefore carry the same underlying abelian group;
  both unwind to nested restriction-of-scalars towers along the structure-sheaf
  global-sections maps. These towers are reconciled by applying the
  restriction-of-scalars composition identity twice and substituting the ring
  equation of
  Lemma~\ref{mk-lem:globalSectionsIso_hom_comp_specMap_appTop}, which shows the
  resulting \(R\)-action is exactly restriction of scalars along \(\varphi\).
\end{proof}

\begin{lemma}[Global sections of an affine pushforward of a tilde-module]\label{mk-lem:gammaPushforwardTildeIso}
  \dcref{custom}

  \uses{mk-lem:gammaPushforwardIso}
  Let \(\varphi : R \to R'\) be a homomorphism of commutative rings and let \(M\)
  be an \(R'\)-module. Then there is a canonical isomorphism of \(R\)-modules
  \[
    \Gamma\bigl((\operatorname{Spec}\varphi)_* \widetilde{M}\bigr)
    \;\cong\;
    \operatorname{restr}_\varphi M .
  \]
\end{lemma}

\begin{proof}
  \uses{mk-lem:gammaPushforwardIso}
  Specialise Lemma~\ref{mk-lem:gammaPushforwardIso} to \(N = \widetilde{M}\) and
  compose with the tilde-unit isomorphism
  \(\Gamma(\widetilde{M}) \cong M\) of \(R'\)-modules (which restricts to an
  isomorphism of \(R\)-modules after applying \(\operatorname{restr}_\varphi\)).
\end{proof}

\begin{lemma}[Sections of an affine pushforward over an arbitrary open]\label{mk-lem:gammaPushforwardIsoAt}
  \dcref{custom}

  \uses{mk-lem:globalSectionsIso_hom_comp_specMap_appTop}
  Let \(\varphi : R \to R'\) be a homomorphism of commutative rings, let \(N\) be
  an \(\operatorname{Spec}(R')\)-module, and let \(U\) be an open of
  \(\operatorname{Spec}(R)\), with preimage
  \(V = (\operatorname{Spec}\varphi)^{-1}(U)\) in \(\operatorname{Spec}(R')\). Then
  there is a canonical isomorphism of \(R\)-modules
  \[
    \Gamma\bigl((\operatorname{Spec}\varphi)_* N,\, U\bigr)
    \;\cong\;
    \operatorname{restr}_\varphi\bigl(\Gamma(N,\, V)\bigr) .
  \]
  This is the open-indexed generalization of
  Lemma~\ref{mk-lem:gammaPushforwardIso}, which is the case \(U = \top\) (so that
  \(V = \top\) as well). The family \(\{e_U\}_U\) of these isomorphisms is moreover
  natural in the open argument: the structure-sheaf restriction maps on the two
  sides commute with the isomorphisms, so the family is a natural isomorphism of the
  two presheaves of \(R\)-modules in the open variable (this naturality is
  established in the proof below). For \(U = D(a)\) it is exactly the linear
  equivalence \(e_{D(a)}\) of movement~(1) in the proof of
  Lemma~\ref{mk-lem:pushforward_spec_tilde_iso}.
\end{lemma}

\begin{proof}
  \uses{mk-lem:globalSectionsIso_hom_comp_specMap_appTop}
  The construction copies that of Lemma~\ref{mk-lem:gammaPushforwardIso} verbatim,
  with the top open \(\top\) replaced throughout by the arbitrary open \(U\) (and
  its preimage \(V\)). The point that makes this replacement free is that the
  \(R\)-module structure on the sections is supplied \emph{uniformly} by
  restriction of scalars along the global-sections ring map at the top open,
  independently of the evaluation open; consequently both sides unwind to the same
  nested restriction-of-scalars towers, reconciled by applying the
  restriction-of-scalars composition identity twice and substituting the
  \(\top\)-level ring equation of
  Lemma~\ref{mk-lem:globalSectionsIso_hom_comp_specMap_appTop}. Naturality in the open
  follows because every constituent of the construction is either a structure-sheaf
  restriction map or conjugation by a fixed ring map, and all of these commute with
  further restriction along an inclusion of opens.

  \medskip
  \emph{Naturality of the family \(\{e_U\}_U\) in the open (remark).} The family is
  in fact a natural isomorphism between two presheaves of \(R\)-modules on
  \(\operatorname{Spec}(R)\): the source presheaf
  \(U \mapsto \Gamma\bigl((\operatorname{Spec}\varphi)_* N,\, U\bigr)\), with the
  structure-sheaf restriction maps of the pushforward, and the target presheaf
  \(U \mapsto
    \operatorname{restr}_\varphi\bigl(\Gamma(N,\, (\operatorname{Spec}\varphi)^{-1}U)\bigr)\),
  with the structure-sheaf restriction maps of \(N\) transported through the preimage
  operation \((\operatorname{Spec}\varphi)^{-1}\) and read as \(R\)-linear maps by
  restriction of scalars along \(\varphi\). Concretely, for any inclusion of opens
  \(U' \subseteq U\), with \(V = (\operatorname{Spec}\varphi)^{-1}U\) and
  \(V' = (\operatorname{Spec}\varphi)^{-1}U'\), the square with horizontal maps
  \(e_U, e_{U'}\) and vertical maps the respective structure-sheaf restrictions
  commutes. The reason is structural: each \(e_U\) is a composite whose constituents
  are of exactly two kinds — structure-sheaf restriction maps (of \(N\),
  respectively of \((\operatorname{Spec}\varphi)_* N\)) and identity-on-carrier
  \(\operatorname{restrictScalars}\) repackagings (conjugation by the \emph{fixed},
  open-independent ring map supplied by the restriction-of-scalars reconciliation and
  by the \(\top\)-level ring equation of
  Lemma~\ref{mk-lem:globalSectionsIso_hom_comp_specMap_appTop}). Restriction maps commute
  with further restriction by the presheaf functoriality axiom (using
  \(V' = (\operatorname{Spec}\varphi)^{-1}U' \subseteq
  (\operatorname{Spec}\varphi)^{-1}U = V\), as \((\operatorname{Spec}\varphi)^{-1}\) is
  monotone), and conjugation by a fixed ring map commutes with restriction because the
  ring map does not depend on \(U\); the naturality square is the paste of these
  component squares, and each \(e_U\) being an isomorphism makes the family a natural
  isomorphism. This open-naturality is exactly the compatibility invoked, at the
  inclusion \(D(a) \subseteq \top\), in the proof of
  Lemma~\ref{mk-lem:pushforward_spec_tilde_iso}.
\end{proof}

\begin{lemma}[The tilde restriction from the top open to a basic open is a localization]\label{mk-lem:tildeRestriction_isLocalizedModule}
  \dcref{custom}

  Let \(M\) be an \(R'\)-module and \(b \in R'\). Then the structure-sheaf
  restriction map
  \[
    \Gamma(\widetilde M,\, \top) \longrightarrow \Gamma(\widetilde M,\, D(b)),
  \]
  read as a map of \(R'\)-modules, exhibits \(\Gamma(\widetilde M, D(b))\) as the
  localization of \(\Gamma(\widetilde M, \top)\) at \(\mathrm{powers}(b)\). This is
  the \(R'\)-side localization fact transported in movement~(3) of the proof of
  Lemma~\ref{mk-lem:pushforward_spec_tilde_iso}.
\end{lemma}

\begin{proof}

By the defining property of \(\widetilde{(-)}\), the tilde-localization map
  \(M \to \Gamma(\widetilde M, D(b))\) exhibits its target as the localization of
  \(M\) at \(\mathrm{powers}(b)\). The global-sections map
  \(M \to \Gamma(\widetilde M, \top)\) is the localization of \(M\) at the trivial
  submonoid \(\mathrm{powers}(1)\) — since \(D(1) = \top\) — and is therefore
  bijective. The triangle identity relating these two localization maps to the
  structure-sheaf restriction \(\Gamma(\widetilde M, \top) \to
  \Gamma(\widetilde M, D(b))\) then exhibits the restriction itself, post-composed
  with the inverse of the bijection \(M \cong \Gamma(\widetilde M, \top)\), as the
  localization of \(\Gamma(\widetilde M, \top)\) at \(\mathrm{powers}(b)\).
\end{proof}

We record three further bespoke supports that drive the basic-open verification of
the affine-pushforward isomorphism below. They carry no external citation and stand
on their proof sketches.

\begin{lemma}[Restriction of scalars of a localized module]\label{mk-lem:powers_restrictScalars}
  \dcref{custom}

  Let \(A\) be a commutative \(R\)-algebra, \(S \subseteq R\) a submonoid, and
  \(f : M \to N\) an \(A\)-linear map that exhibits \(N\) as the localization of
  \(M\) at the image submonoid
  \(\operatorname{algMap}_A(S) \subseteq A\). Then the underlying \(R\)-linear map
  \(\operatorname{restr}_R f\) exhibits \(N\) as the localization of \(M\) at \(S\)
  itself.
\end{lemma}

\begin{proof}


This is the exact converse of the standard fact that a localization of \(M\) at a
  submonoid \(S\) of \(R\) remains a localization after enlarging the base to an
  \(R\)-algebra \(A\) along the image submonoid \(\operatorname{algMap}_A(S)\). One
  checks the three defining conditions of a localized module directly. Each
  \(s \in S\) acts invertibly on \(N\) because its image
  \(\operatorname{algMap}_A(s)\) does and the two actions coincide by the scalar
  tower \(R \to A \to N\). Surjectivity of \(f\) up to denominators and the
  equality-witness condition transport along the map
  \(s \mapsto \operatorname{algMap}_A(s)\) from \(S\) onto
  \(\operatorname{algMap}_A(S)\), again using
  \(\operatorname{algMap}_A(s)\cdot{} = s\cdot{}\). In the application
  \(A = R'\), \(S = \mathrm{powers}(a)\), one has
  \(\operatorname{algMap}_{R'}(\mathrm{powers}(a)) = \mathrm{powers}(\varphi a)\), so
  this identifies the localization of \(\operatorname{restr}_\varphi M\) at \(a\)
  with the localization of \(M\) at \(\varphi a\); it is the ring-change core of the
  basic-open computation below.
\end{proof}

\begin{lemma}[Section-level counit isomorphism from a localization hypothesis]\label{mk-lem:fromTildeGamma_app_isIso_of_localized}
  \dcref{custom}

  Let \(N\) be an \(\operatorname{Spec}(R)\)-module and \(a \in R\). Suppose the
  structure-sheaf restriction map \(\Gamma(N, \top) \to \Gamma(N, D(a))\), read as
  an \(R\)-linear map of the global-section modules, exhibits \(\Gamma(N, D(a))\) as
  the localization of \(\Gamma(N, \top)\) at \(\mathrm{powers}(a)\). Then the
  tilde--\(\Gamma\) counit
  \(\mathrm{fromTilde}\Gamma\,N : \widetilde{\Gamma(N,\top)} \to N\) restricts to an
  isomorphism on the sections over the basic open \(D(a)\).
\end{lemma}

\begin{proof}


Over \(D(a)\) the section map \(L\) of the counit sits in the triangle identity
  \(L \circ j = \rho\), where
  \(j : \Gamma(N,\top) \to \Gamma\bigl(\widetilde{\Gamma(N,\top)}, D(a)\bigr)\) is
  the tilde-localization map and \(\rho : \Gamma(N,\top) \to \Gamma(N, D(a))\) is
  the structure-sheaf restriction. Both \(j\) and \(\rho\) exhibit their targets as
  the localization of \(\Gamma(N,\top)\) at \(\mathrm{powers}(a)\) — \(j\) by the
  defining property of \(\widetilde{(-)}\), and \(\rho\) by hypothesis. By the
  uniqueness of localized modules there is a unique \(R\)-linear isomorphism \(e\)
  with \(e \circ j = \rho\); the triangle identity then forces \(L = e\), so \(L\) is
  an isomorphism. (This is the per-basic-open input feeding the basis-local
  criterion Lemma~\ref{mk-lem:modules_isIso_of_isBasis}.)
\end{proof}

\begin{lemma}[Affine pushforward of a tilde-module, conditional form]\label{mk-lem:pushforward_spec_tilde_iso_conditional}
  \dcref{custom}

  \uses{mk-lem:modules_isIso_of_isBasis, mk-lem:gammaPushforwardTildeIso,
  mk-lem:fromTildeGamma_app_isIso_of_localized}
  Let \(\varphi : R \to R'\) be a homomorphism of commutative rings and \(M\) an
  \(R'\)-module. Suppose that for every \(a \in R\) the structure-sheaf restriction
  of the pushforward \(N := (\operatorname{Spec}\varphi)_*\widetilde M\) exhibits its
  sections over \(D(a)\) as the localization of its global sections at
  \(\mathrm{powers}(a)\). Then there is a canonical isomorphism of
  \(\operatorname{Spec}(R)\)-modules
  \[
    (\operatorname{Spec}\varphi)_*\widetilde M \;\cong\;
    \widetilde{\bigl(\operatorname{restr}_\varphi M\bigr)} .
  \]
\end{lemma}

\begin{proof}
  \uses{mk-lem:modules_isIso_of_isBasis, mk-lem:gammaPushforwardTildeIso, mk-lem:fromTildeGamma_app_isIso_of_localized}
  Write \(N := (\operatorname{Spec}\varphi)_*\widetilde M\). Under the hypothesis,
  Lemma~\ref{mk-lem:fromTildeGamma_app_isIso_of_localized} applies for each \(a \in R\),
  so the counit \(\mathrm{fromTilde}\Gamma\,N\) restricts to an isomorphism on the
  sections over every basic open \(D(a)\). Since the basic opens form a basis of
  \(\operatorname{Spec}(R)\), the basis-local criterion
  Lemma~\ref{mk-lem:modules_isIso_of_isBasis} upgrades this to: the counit
  \(\mathrm{fromTilde}\Gamma\,N\) is an isomorphism, i.e.\ \(N\) lies in the
  essential image of \(\widetilde{(-)}\). Composing the inverse counit with the
  tilde of the global-sections comparison
  Lemma~\ref{mk-lem:gammaPushforwardTildeIso} yields the asserted object isomorphism.
\end{proof}

\begin{lemma}[Affine pushforward of a tilde-module]\label{mk-lem:pushforward_spec_tilde_iso}
  \dcref{custom}

  \uses{mk-lem:modules_isIso_of_isBasis, mk-lem:gammaPushforwardTildeIso,
        mk-lem:pushforward_spec_tilde_iso_conditional, mk-lem:powers_restrictScalars,
        mk-lem:gammaPushforwardIso, mk-lem:gammaPushforwardIsoAt,
        mk-lem:tildeRestriction_isLocalizedModule}
  Let \(\varphi : R \to R'\) be a homomorphism of commutative rings and let \(M\)
  be an \(R'\)-module. Then there is a canonical isomorphism of
  \(\operatorname{Spec}(R)\)-modules
  \[
    \bigl(\operatorname{Spec}\varphi\bigr)_*\,\widetilde{M}
    \;\cong\;
    \widetilde{\bigl(\operatorname{restr}_\varphi M\bigr)},
  \]
  where \(\operatorname{Spec}\varphi : \operatorname{Spec}(R') \to
  \operatorname{Spec}(R)\) is the induced morphism of affine schemes,
  \(\widetilde{(-)}\) denotes the quasi-coherent sheaf associated to a module, and
  \(\operatorname{restr}_\varphi M\) is the \(R\)-module obtained from \(M\) by
  restriction of scalars along \(\varphi\). That is, pushing the quasi-coherent
  sheaf \(\widetilde{M}\) forward along \(\operatorname{Spec}\varphi\) is, up to
  this canonical identification, the tilde of the restriction of scalars of \(M\).

  This is the classical fact that the direct image of a quasi-coherent sheaf along
  a morphism of affine schemes is computed on modules by restriction of scalars; it
  is the affine companion of the pullback formula (the affine case of the
  pushforward analogue of Stacks ``$\widetilde{M}$ and pullback''). It is recorded
  here as bespoke project infrastructure because the affine-pushforward
  identification does not appear as a ready-made result in the existing
  quasi-coherent sheaf theory.
\end{lemma}

\begin{proof}
  \uses{mk-lem:modules_isIso_of_isBasis, mk-lem:gammaPushforwardTildeIso, mk-lem:pushforward_spec_tilde_iso_conditional, mk-lem:powers_restrictScalars, mk-lem:gammaPushforwardIso, mk-lem:gammaPushforwardIsoAt, mk-lem:tildeRestriction_isLocalizedModule}
  We build the isomorphism \emph{directly on a basis of basic opens}. This is the
  non-circular route: checking the comparison map on basic opens yields the object
  isomorphism, and quasi-coherence of the pushforward is then a corollary, not a
  prerequisite. (One must not reverse this order: the global-sections comparison
  \(\Gamma\bigl((\operatorname{Spec}\varphi)_*\widetilde{M}\bigr) \cong
  \operatorname{restr}_\varphi M\) of
  Lemma~\ref{mk-lem:gammaPushforwardTildeIso} alone would only upgrade to an object
  isomorphism once the pushforward is already known to be quasi-coherent, so
  quasi-coherence cannot be deduced from an object isomorphism built that way.)

  \emph{The comparison morphism.} There is a canonical morphism of
  \(\operatorname{Spec}(R)\)-modules
  \[
    \alpha : \widetilde{\bigl(\operatorname{restr}_\varphi M\bigr)}
    \longrightarrow \bigl(\operatorname{Spec}\varphi\bigr)_*\widetilde{M},
  \]
  namely the tilde-adjoint of the global-sections identification of
  Lemma~\ref{mk-lem:gammaPushforwardTildeIso}. We show \(\alpha\) is an isomorphism.

  \medskip
  \emph{Reduction to basic opens.} The basic opens \(D(a)\), for \(a \in R\), form a
  basis of the topology of \(\operatorname{Spec}(R)\). By the basis-local criterion
  Lemma~\ref{mk-lem:modules_isIso_of_isBasis} it suffices to check that \(\alpha\)
  restricts to an isomorphism on the sections over each basic open \(D(a)\).

  \medskip
  \emph{The computation over \(D(a)\).} On the source, the sections of
  \(\widetilde{(\operatorname{restr}_\varphi M)}\) over \(D(a)\) are the
  localization \(\bigl(\operatorname{restr}_\varphi M\bigr)[1/a]\) as an
  \(R[1/a]\)-module. On the target, since \((\operatorname{Spec}\varphi)^{-1}D(a) =
  D(\varphi a)\), the sections of \(\bigl(\operatorname{Spec}\varphi\bigr)_*
  \widetilde{M}\) over \(D(a)\) are the sections of \(\widetilde{M}\) over
  \(D(\varphi a)\), i.e.\ the localization \(M[1/\varphi a]\), regarded as an
  \(R[1/a]\)-module by restriction of scalars along the induced map
  \(R[1/a] \to R'[1/\varphi a]\). These two localizations agree: the element \(a\)
  acts on \(\operatorname{restr}_\varphi M\) through \(\varphi a\), so localizing
  \(\operatorname{restr}_\varphi M\) at the multiplicative set generated by \(a\) is
  the same module as localizing \(M\) at the multiplicative set generated by
  \(\varphi a\). Both are the universal localization of
  \(M\) inverting \(\varphi a\), so the comparison map \(\alpha\) is an isomorphism
  on the sections over \(D(a)\). As this holds for every basic open, \(\alpha\) is
  an isomorphism.

  \medskip
  \emph{The explicit affine route over \(D(a)\).} The agreement of the two
  localizations just described is the entire content, but it must be realized
  through the global ring \(R\), because the \(R\)-action on the pushforward
  sections over \(D(a)\) is built from the global-sections module of
  \(N := (\operatorname{Spec}\varphi)_*\widetilde M\) rather than directly from
  \(R[1/a]\). The reduction therefore proceeds exactly as the conditional
  assembly Lemma~\ref{mk-lem:pushforward_spec_tilde_iso_conditional} prescribes: it
  suffices to supply, for each \(a \in R\), the hypothesis \(\mathrm{hloc}(a)\) that
  the structure-sheaf restriction \(\Gamma(N, \top) \to \Gamma(N, D(a))\) exhibits
  \(\Gamma(N, D(a))\) as the localization of \(\Gamma(N, \top)\) at
  \(\mathrm{powers}(a)\). The \(R\)-action on \(\Gamma(N, D(a))\) is the honest
  restriction-of-scalars action carried by the scalar tower \(R \to R' \to {-}\)
  determined by \(\varphi\); against that structure the ring-change converse
  Lemma~\ref{mk-lem:powers_restrictScalars} is applicable, so the carrier of the
  \(R\)-module is exactly the one to which the localization claim refers, and no
  recursively-defined or ad-hoc section-level scalar action intervenes.
  We obtain \(\mathrm{hloc}(a)\) by transporting the already-known \(R'\)-side
  localization (Lemma~\ref{mk-lem:tildeRestriction_isLocalizedModule}), in
  element-free fashion, across the open-indexed section comparison of
  Lemma~\ref{mk-lem:gammaPushforwardIsoAt} \emph{using its naturality in the open}
  (the naturality remark in the proof of
  Lemma~\ref{mk-lem:gammaPushforwardIsoAt}). This proceeds in three
  movements.

  \emph{(1) The \(D(a)\)-level linear equivalence.} This is the component at
  \(U = D(a)\) of the open-indexed family of
  Lemma~\ref{mk-lem:gammaPushforwardIsoAt}: the \(R\)-linear isomorphism
  \[
    e_{D(a)} : \Gamma\bigl((\operatorname{Spec}\varphi)_*\widetilde M,\, D(a)\bigr)
      \;\cong\;
      \operatorname{restr}_\varphi\bigl(\Gamma(\widetilde M,\, D(\varphi a))\bigr) ,
  \]
  using \((\operatorname{Spec}\varphi)^{-1}D(a) = D(\varphi a)\), which is
  definitional. No fresh \(D(a)\)-level construction is needed: this is precisely the
  open-indexed isomorphism already supplied by
  Lemma~\ref{mk-lem:gammaPushforwardIsoAt} evaluated at the basic open \(D(a)\).

  \emph{(2) The naturality square for \(D(a) \subseteq \top\).} The compatibility
  that relates \(e_{D(a)}\) to the \(\top\)-level isomorphism
  \(e_\top\) (the global-sections comparison of
  Lemma~\ref{mk-lem:gammaPushforwardIso}) is precisely the open-naturality square of the
  family \(\{e_U\}_U\) (established in the proof of
  Lemma~\ref{mk-lem:gammaPushforwardIsoAt}), instantiated at the inclusion
  \(D(a) \subseteq \top\): the square
  \[
    \begin{array}{ccc}
      \Gamma\bigl((\operatorname{Spec}\varphi)_*\widetilde M,\, \top\bigr)
        & \xrightarrow{\;e_\top\;}
        & \operatorname{restr}_\varphi\bigl(\Gamma(\widetilde M,\, \top)\bigr) \\[4pt]
      {\scriptstyle \mathrm{res}_{\top \to D(a)}}\big\downarrow & &
        \big\downarrow{\scriptstyle \mathrm{res}_{\top \to D(\varphi a)}} \\[4pt]
      \Gamma\bigl((\operatorname{Spec}\varphi)_*\widetilde M,\, D(a)\bigr)
        & \xrightarrow{\;e_{D(a)}\;}
        & \operatorname{restr}_\varphi\bigl(\Gamma(\widetilde M,\, D(\varphi a))\bigr)
    \end{array}
  \]
  commutes, using \((\operatorname{Spec}\varphi)^{-1}D(a) = D(\varphi a)\), which is
  definitional. Because the family \(\{e_U\}_U\) is a natural isomorphism, this single
  square is available uniformly in \(a\); there is no per-\(a\) section-level identity
  to re-prove by hand. It intertwines the source restriction
  \(\mathrm{res}_{\top \to D(a)}\) on the pushforward with the (restriction-of-scalars
  reading of the) target restriction \(\mathrm{res}_{\top \to D(\varphi a)}\) on
  \(\widetilde M\).

  \emph{(3) Discharging \(\mathrm{hloc}(a)\) by transport.} The target restriction
  \(\mathrm{res}_{\top \to D(\varphi a)} : \Gamma(\widetilde M, \top) \to
  \Gamma(\widetilde M, D(\varphi a))\) exhibits \(\Gamma(\widetilde M, D(\varphi a)) =
  M[1/\varphi a]\) as the \(\mathrm{powers}(\varphi a)\)-localization of
  \(M = \Gamma(\widetilde M, \top)\) over \(R'\); this is exactly
  Lemma~\ref{mk-lem:tildeRestriction_isLocalizedModule}. Apply the ring-change converse
  Lemma~\ref{mk-lem:powers_restrictScalars} (with \(R\)-algebra \(A = R'\) and submonoid
  \(\mathrm{powers}(a)\), so that
  \(\operatorname{algMap}_{R'}(\mathrm{powers}(a)) = \mathrm{powers}(\varphi a)\)) —
  its \(\operatorname{Algebra} R\, R'\) and scalar-tower hypotheses being supplied by
  the restriction-of-scalars structure along \(\varphi\) — to read this same map as
  the \(\mathrm{powers}(a)\)-localization over \(R\). The commuting naturality square
  of movement (2) intertwines this target restriction with the source restriction
  \(\mathrm{res}_{\top \to D(a)} : \Gamma(N, \top) \to \Gamma(N, D(a))\) by way of the
  isomorphisms \(e_\top\) and \(e_{D(a)}\); transporting the localization property
  across that square — via the transport-along-an-equivalence principle for localized
  modules — yields that \(\mathrm{res}_{\top \to D(a)}\) is the
  \(\mathrm{powers}(a)\)-localization on \(\Gamma(N, D(a))\) over \(R\), i.e.\
  \(\mathrm{hloc}(a)\). No section-level computation specific to the individual \(a\)
  is performed: the only \(a\)-dependent input is the instantiation of the single
  natural-isomorphism square at \(D(a)\). Feeding the family
  \(\{\mathrm{hloc}(a)\}_{a \in R}\) to the conditional assembly
  Lemma~\ref{mk-lem:pushforward_spec_tilde_iso_conditional} produces the unconditional
  object isomorphism \(\alpha\).

  \medskip
  \emph{Naturality in the open drives the transport.} The comparison assembled
  here, \((\operatorname{Spec}\varphi)_* \circ \widetilde{(-)} \cong
  \widetilde{(-)} \circ \operatorname{restr}_\varphi\), is natural in \emph{both}
  arguments — in the module \(M\) and in the open \(U\). The naturality in the
  module is the usual functoriality; the naturality in the open is the naturality
  remark in the proof of Lemma~\ref{mk-lem:gammaPushforwardIsoAt}: the family
  \(\{e_U\}_U\) of Lemma~\ref{mk-lem:gammaPushforwardIsoAt} commutes with the
  structure-sheaf restriction maps on both sides. It is precisely this
  open-naturality that drives the localization transport \emph{uniformly}: the
  per-\(a\) localization fact \(\mathrm{hloc}(a)\) follows from the \(\top\)-level
  localization Lemma~\ref{mk-lem:tildeRestriction_isLocalizedModule} together with that
  naturality square for the inclusion
  \(D(a) \hookrightarrow \top\), rather than from a section-level naturality square
  re-proved by hand for each individual \(a\). This is the single remaining
  obligation of the construction; the carrier of the \(R\)-action on the pushforward
  sections is the restriction-of-scalars structure along \(\varphi\), against which
  Lemma~\ref{mk-lem:powers_restrictScalars} applies directly.

  The \(R\)-module structure on the \(R'\)-side sections to which the ring-change
  converse Lemma~\ref{mk-lem:powers_restrictScalars} is applied is the honest
  restriction-of-scalars structure: \(R\) acts through \(\varphi\), i.e.\ via the
  \(R\)-algebra structure on \(R'\) determined by \(\varphi\) and the associated
  scalar tower \(R \to R' \to M[1/\varphi a]\). It is \emph{not} an ad-hoc,
  recursively-defined scalar action. Lemma~\ref{mk-lem:powers_restrictScalars} is
  stated against exactly this restriction-of-scalars / scalar-tower structure, and
  it is under that structure that localizing \(\operatorname{restr}_\varphi M\) at
  \(a\) coincides with localizing \(M\) at \(\varphi a\).
  \medskip
  \emph{Quasi-coherence as a corollary.} The object isomorphism just constructed
  exhibits \(\bigl(\operatorname{Spec}\varphi\bigr)_*\widetilde{M} \cong
  \widetilde{(\operatorname{restr}_\varphi M)}\) as the tilde of a module. Tilde
  modules are quasi-coherent and quasi-coherence of \(\operatorname{Spec}(R)\)-modules
  is closed under isomorphism, so the pushforward
  \(\bigl(\operatorname{Spec}\varphi\bigr)_*\widetilde{M}\) is quasi-coherent. This
  conclusion is stated \emph{after} the object isomorphism, of which it is a
  consequence — no separate ``pushforward preserves quasi-coherent'' input is
  needed for the affine lane.
\end{proof}

\begin{remark}
  This single isomorphism discharges both inputs that the affine reduction below
  requires of the pushforward of a tilde-module.
  \begin{enumerate}
    \item \emph{The global-sections comparison.} Applying the global-sections
      (module-of-sections) functor to the isomorphism identifies the
      \(R\)-module of sections of
      \(\bigl(\operatorname{Spec}\varphi\bigr)_*\widetilde{M}\) with
      \(\operatorname{restr}_\varphi M\); this is exactly the \(\Gamma\)-fragment
      comparison used to recognise the section-level base-change map.
    \item \emph{Quasi-coherence of the pushforward.} The tilde of any module is
      quasi-coherent, and quasi-coherence of \(\operatorname{Spec}(R)\)-modules is
      closed under isomorphism; hence
      \(\bigl(\operatorname{Spec}\varphi\bigr)_*\widetilde{M}\) is quasi-coherent.
      In particular no separate ``pushforward preserves quasi-coherent'' theorem is
      needed for the affine lane.
  \end{enumerate}
  Because both inputs the affine reduction asks of the pushforward are corollaries
  of this single object isomorphism — the section-level module identification and
  the quasi-coherence of the pushforward — this lemma is the sole bespoke
  ingredient the affine lane requires.
\end{remark}

\subsection{The pullback companion}

The affine reduction below requires, alongside the pushforward dictionary, the
\emph{pullback} companion: the affine description of the inverse image of a
tilde-module. It is the affine case of the classical fact that the pullback of a
quasi-coherent sheaf is computed on modules by base change of scalars.

The pullback companion is built by uniqueness of left adjoints, which is driven by
the following packaging of the per-object \(\Gamma\)-fragment comparison as a
natural isomorphism of right-adjoint functors. It is bespoke project
infrastructure and carries no external citation.

\begin{lemma}[The pushforward--global-sections natural isomorphism]\label{mk-lem:gammaPushforwardNatIso}
  \dcref{custom}

  \uses{mk-lem:gammaPushforwardIso}
  Let \(\varphi : R \to R'\) be a homomorphism of commutative rings. Then the
  per-object global-sections comparison of
  Lemma~\ref{mk-lem:gammaPushforwardIso} assembles into a natural isomorphism of
  functors \((\operatorname{Spec} R').\mathrm{Modules} \to \operatorname{ModuleCat}(R)\)
  \[
    \bigl(\operatorname{Spec}\varphi\bigr)_* \;\circ\; \Gamma_R
    \;\cong\;
    \Gamma_{R'} \;\circ\; \operatorname{restr}_\varphi ,
  \]
  whose component at an \(\operatorname{Spec}(R')\)-module \(N\) is the isomorphism
  \(\Gamma\bigl((\operatorname{Spec}\varphi)_* N\bigr) \cong
  \operatorname{restr}_\varphi(\Gamma N)\) of
  Lemma~\ref{mk-lem:gammaPushforwardIso}.
\end{lemma}

\begin{proof}
  \uses{mk-lem:gammaPushforwardIso}
  The components are the isomorphisms supplied by
  Lemma~\ref{mk-lem:gammaPushforwardIso}. Naturality in the module argument \(N\)
  holds on the nose: every constituent of the comparison
  (the restriction-of-scalars composition repackagings and the
  restriction-of-scalars congruence) is the identity on underlying elements,
  merely re-presenting the module structure on an unchanged carrier, so each
  naturality square is a pointwise identity. This is the right-adjoint natural
  isomorphism that drives the uniqueness-of-left-adjoints argument for the affine
  pullback dictionary Lemma~\ref{mk-lem:pullback_spec_tilde_iso}: identifying the two
  right adjoints \((\operatorname{Spec}\varphi)_* \circ \Gamma_R\) and
  \(\Gamma_{R'} \circ \operatorname{restr}_\varphi\) by this isomorphism forces the
  corresponding left adjoints to be isomorphic.
\end{proof}

\begin{lemma}[Affine pullback of a tilde-module]\label{mk-lem:pullback_spec_tilde_iso}
  \dcref{01I9,custom}

  \uses{mk-lem:pushforward_spec_tilde_iso}
  \textit{Source: Stacks Project, Schemes, Lemma ``\(\widetilde M\) and pullback''
  (Tag 01I9), part (1).}
  Let \(\varphi : R \to R'\) be a homomorphism of commutative rings and let \(M\)
  be an \(R\)-module. Then there is a canonical isomorphism of
  \(\operatorname{Spec}(R')\)-modules
  \[
    \bigl(\operatorname{Spec}\varphi\bigr)^{*}\,\widetilde{M}
    \;\cong\;
    \widetilde{\bigl(R' \otimes_R M\bigr)},
  \]
  where \(\operatorname{Spec}\varphi : \operatorname{Spec}(R') \to
  \operatorname{Spec}(R)\) is the induced morphism of affine schemes,
  \(\widetilde{(-)}\) denotes the quasi-coherent sheaf associated to a module, and
  \(R' \otimes_R M\) is the base change of \(M\) along \(\varphi\) — the
  \(R'\)-module obtained by extension of scalars. The isomorphism is natural in
  \(M\). That is, pulling the quasi-coherent sheaf \(\widetilde{M}\) back along
  \(\operatorname{Spec}\varphi\) is, up to this canonical identification, the tilde
  of the base change of \(M\). It is the pullback companion of the affine
  pushforward dictionary Lemma~\ref{mk-lem:pushforward_spec_tilde_iso}, and is part~(1)
  of the same Stacks lemma of which the pushforward formula is part~(2).
\end{lemma}
\begin{proof}
  \uses{mk-lem:pushforward_spec_tilde_iso}
  Pullback along \(\operatorname{Spec}\varphi\) is, on quasi-coherent sheaves, the
  extension-of-scalars operation \(- \otimes_R R'\) on modules transported through
  the \(\widetilde{(-)}\) / \(\operatorname{moduleSpec}\Gamma\) equivalence between
  modules and quasi-coherent sheaves. Concretely, the global sections of
  \(\bigl(\operatorname{Spec}\varphi\bigr)^{*}\widetilde{M}\) are the base change
  \(R' \otimes_R M\): the pullback functor \((\operatorname{Spec}\varphi)^{*}\) is
  left adjoint to the pushforward \((\operatorname{Spec}\varphi)_{*}\), and under
  the global-sections identification the latter is restriction of scalars along
  \(\varphi\) (Lemma~\ref{mk-lem:pushforward_spec_tilde_iso}); the left adjoint of
  restriction of scalars is extension of scalars \(M \mapsto R' \otimes_R M\). The
  comparison map is the unit of the base-change adjunction, and the universal
  property of base change identifies it with the canonical isomorphism
  \(\bigl(\operatorname{Spec}\varphi\bigr)^{*}\widetilde{M} \cong
  \widetilde{(R' \otimes_R M)}\). Naturality in \(M\) is the naturality of that
  unit. (One may equally read this off the standard-open description: on
  \(D(\varphi a)\) the localization \((R' \otimes_R M)[1/a]\) is
  \(R'[1/\varphi a] \otimes_{R[1/a]} M[1/a]\), which is the base change of the
  sections of \(\widetilde M\) over \(D(a)\).)
\end{proof}

\subsection{The concrete-tilde affine base-change brick (active route)}
\label{mk-subsec:fbc_concrete_brick}

The key affine input for this construction is the following isomorphism. It
is assembled entirely from the two tilde dictionaries
(Lemmas~\ref{mk-lem:pushforward_spec_tilde_iso} and~\ref{mk-lem:pullback_spec_tilde_iso}) and
Mathlib's cancellation isomorphism, and --- crucially --- it never forms the adjoint
mate of Definition~\ref{mk-def:pushforward_base_change_map}, so it does not re-incur the
mate\(\,\leftrightarrow\,\operatorname{cancelBaseChange}\) obligation of the
superseded canonical route.

\begin{lemma}[Cancellation isomorphism for iterated base change]
  \label{mk-lem:cancelBaseChange_mathlib}
  \dcref{custom}

  Let \(R \to R'\) and \(R \to A\) be commutative-ring homomorphisms and \(M\) an
  \(A\)-module. There is a canonical isomorphism of \(R'\)-modules
  \[
    \operatorname{cancelBaseChange} :
    (R' \otimes_R A) \otimes_A M \;\xrightarrow{\ \sim\ }\; R' \otimes_R M,
    \qquad (r' \otimes a) \otimes m \;\mapsto\; r' \otimes (a \cdot m),
  \]
  with no flatness hypothesis whatsoever.
\end{lemma}

\begin{lemma}[Module-level base-change cancellation iso]\label{mk-lem:baseChangeCancelModuleIso}
  \dcref{custom}

  \uses{mk-lem:cancelBaseChange_mathlib}
  Let \(\varphi : R \to A\), \(\psi : R \to R'\), \(\rho : A \to B\),
  \(\sigma : R' \to B\) be commutative-ring homomorphisms forming a pushout square in
  \(\mathrm{CommRing}\) (so that \(B\) is, up to canonical isomorphism, the tensor
  product \(A \otimes_R R'\)), and let \(M\) be an \(A\)-module. Then there is a
  canonical isomorphism of \(R'\)-modules
  \[
    \operatorname{extendScalars}_\psi\bigl(\operatorname{restr}_\varphi M\bigr)
    \;\cong\;
    \operatorname{restr}_\sigma\bigl(\operatorname{extendScalars}_\rho M\bigr),
  \]
  i.e.\ \(R' \otimes_R M \cong \operatorname{restr}_\sigma\bigl(B \otimes_A M\bigr)\),
  realising the inverse \(\operatorname{cancelBaseChange}^{-1}\) of
  Lemma~\ref{mk-lem:cancelBaseChange_mathlib} as a map between the extension-of-scalars
  objects attached to the two legs of the square.
\end{lemma}

\begin{proof}
  \uses{mk-lem:cancelBaseChange_mathlib}
  The pushout hypothesis exhibits \((R, R', A, B)\) as a commutative-algebra pushout.
  For such a square the cancellation isomorphism Lemma~\ref{mk-lem:cancelBaseChange_mathlib} is
  exactly the \(R'\)-linear isomorphism \(B \otimes_A M \cong R' \otimes_R M\); its
  inverse, read through the definitional identifications
  \(R' \otimes_R M = \operatorname{extendScalars}_\psi(\operatorname{restr}_\varphi M)\)
  and \(B \otimes_A M = \operatorname{extendScalars}_\rho M\), is the displayed map. It
  is an isomorphism with no flatness hypothesis.
\end{proof}

\begin{lemma}[Affine termwise base change, abstract-\(B\) framing]\label{mk-lem:affine_pushforward_pullback_baseChange}
  \dcref{02KG,custom}

  \uses{mk-lem:pushforward_spec_tilde_iso, mk-lem:pullback_spec_tilde_iso,
        mk-lem:cancelBaseChange_mathlib, mk-lem:baseChangeCancelModuleIso}
  \textit{Source: Stacks Project, Cohomology of Schemes, Tag 02KG
  (\texttt{lemma-affine-base-change}).}
  Let \(\varphi : R \to A\), \(\psi : R \to R'\), \(\rho : A \to B\),
  \(\sigma : R' \to B\) be commutative-ring homomorphisms forming a pushout square in
  \(\mathrm{CommRing}\),
  \[
    \begin{array}{ccc}
      R & \xrightarrow{\;\varphi\;} & A \\[2pt]
      {\scriptstyle \psi}\big\downarrow & & \big\downarrow{\scriptstyle \rho} \\[2pt]
      R' & \xrightarrow{\;\sigma\;} & B
    \end{array}
  \]
  equivalently --- applying \(\operatorname{Spec}\), which carries a pushout of rings to
  a cartesian (pullback) square of affine schemes --- the cartesian square
  \[
    \begin{array}{ccc}
      \operatorname{Spec} B & \xrightarrow{\;\operatorname{Spec}\rho\;}
        & \operatorname{Spec} A \\[2pt]
      {\scriptstyle \operatorname{Spec}\sigma}\big\downarrow & &
        \big\downarrow{\scriptstyle \operatorname{Spec}\varphi} \\[2pt]
      \operatorname{Spec} R' & \xrightarrow{\;\operatorname{Spec}\psi\;}
        & \operatorname{Spec} R .
    \end{array}
  \]
  Write \(S = \operatorname{Spec} R\), \(S' = \operatorname{Spec} R'\),
  \(X = \operatorname{Spec} A\), \(X' = \operatorname{Spec} B\), with
  \(f = \operatorname{Spec}\varphi\), \(g = \operatorname{Spec}\psi\),
  \(f' = \operatorname{Spec}\sigma\), \(g' = \operatorname{Spec}\rho\); let \(M\) be an
  \(A\)-module and \(\mathcal{F} = \widetilde{M}\). Then there is an isomorphism of
  \(\operatorname{Spec}(R')\)-modules
  \[
    g^*\bigl(f_*\widetilde{M}\bigr)
    \;\cong\;
    f'_*\bigl((g')^*\widetilde{M}\bigr),
  \]
  natural in the module \(M\), built directly from the tilde primitives and the
  cancellation isomorphism, and \emph{never} identified with the canonical
  adjoint-mate base-change map.

\end{lemma}

\begin{proof}
  \uses{mk-lem:pushforward_spec_tilde_iso, mk-lem:pullback_spec_tilde_iso,
        mk-lem:cancelBaseChange_mathlib, mk-lem:baseChangeCancelModuleIso}
        The two sides are transported to modules through the affine dictionaries and then
  compared by the cancellation isomorphism, in three composable moves; each move is an
  isomorphism, so the composite is.

  \emph{(1) The source \(g^*(f_*\widetilde{M})\).} The pushforward dictionary
  Lemma~\ref{mk-lem:pushforward_spec_tilde_iso} identifies \(f_*\widetilde{M}\) over
  \(\operatorname{Spec}(R)\) with \(\widetilde{\operatorname{restr}_\varphi M}\). The
  pullback dictionary Lemma~\ref{mk-lem:pullback_spec_tilde_iso} then identifies its
  pullback \(g^*\bigl(\widetilde{\operatorname{restr}_\varphi M}\bigr)\) with
  \(\widetilde{\bigl(R' \otimes_R M\bigr)}\), the tilde of the \(R'\)-module
  \(\operatorname{extendScalars}_\psi(\operatorname{restr}_\varphi M) = R' \otimes_R M\).

  \emph{(2) The target \(f'_*((g')^*\widetilde{M})\).} The pullback dictionary
  Lemma~\ref{mk-lem:pullback_spec_tilde_iso} (applied to \(g' = \operatorname{Spec}\rho\))
  identifies \((g')^*\widetilde{M}\) with
  \(\widetilde{\bigl(\operatorname{extendScalars}_\rho M\bigr)}
  = \widetilde{\bigl(B \otimes_A M\bigr)}\). The pushforward dictionary
  Lemma~\ref{mk-lem:pushforward_spec_tilde_iso} (applied to
  \(f' = \operatorname{Spec}\sigma\)) identifies the pushforward of this tilde with
  \(\widetilde{\bigl(\operatorname{restr}_\sigma(B \otimes_A M)\bigr)}\), the tilde of
  \(B \otimes_A M\) read as an \(R'\)-module by restriction of scalars along \(\sigma\).

  \emph{(3) Cancellation.} The module-level core
  Lemma~\ref{mk-lem:baseChangeCancelModuleIso} supplies the \(R'\)-linear isomorphism
  \(\operatorname{extendScalars}_\psi(\operatorname{restr}_\varphi M) \cong
  \operatorname{restr}_\sigma(\operatorname{extendScalars}_\rho M)\), i.e.\
  \(R' \otimes_R M \cong \operatorname{restr}_\sigma(B \otimes_A M)\); it is the
  realization of \(\operatorname{cancelBaseChange}^{-1}\)
  (Lemma~\ref{mk-lem:cancelBaseChange_mathlib}) for the square. Applying
  \(\widetilde{(-)}\) and composing the inverse of (1), this iso, and (2) yields the
  asserted isomorphism \(g^*(f_*\widetilde{M}) \cong f'_*((g')^*\widetilde{M})\).

  \medskip
  \emph{Realization as a chain of isomorphisms.} The three moves are assembled as a
  five-step chain of isomorphisms of \(\operatorname{Spec}(R')\)-modules
  \[
    g^*\bigl(f_*\widetilde{M}\bigr)
    \;\xrightarrow{\;\sim\;}\; \widetilde{\bigl(R' \otimes_R M\bigr)}
    \;\xrightarrow{\;\widetilde{(\cdot)}\,\mathrm{of\ core}\;}\;
    \widetilde{\bigl(\operatorname{restr}_\sigma(B \otimes_A M)\bigr)}
    \;\xrightarrow{\;\sim\;} f'_*\bigl((g')^*\widetilde{M}\bigr),
  \]
  in which the first arrow is move (1) (the pushforward dictionary
  Lemma~\ref{mk-lem:pushforward_spec_tilde_iso} followed by the pullback dictionary
  Lemma~\ref{mk-lem:pullback_spec_tilde_iso}), the middle arrow is \(\widetilde{(-)}\)
  applied to the module-level core Lemma~\ref{mk-lem:baseChangeCancelModuleIso}, and the
  last arrow is the inverse of move (2) (the pullback then pushforward dictionaries for
  \(g'\) and \(f'\)). The only step carrying mathematical content beyond the tilde
  dictionaries is the middle one: keeping \(B\) abstract ensures the cancellation
  Lemma~\ref{mk-lem:baseChangeCancelModuleIso} applies from the pushout hypothesis
  directly, with no adjoint mate formed.

  Naturality in \(M\) is the naturality of each constituent: the two tilde dictionaries
  are natural in the module argument, and \(\operatorname{cancelBaseChange}\) (hence the
  module core Lemma~\ref{mk-lem:baseChangeCancelModuleIso}) is natural in \(M\) by Mathlib.
  No adjoint mate is formed at any point, so the
  mate\(\,\leftrightarrow\,\operatorname{cancelBaseChange}\) coherence obligation of the
  canonical route never arises.
\end{proof}

\section{Gluing the affine base-change isomorphisms to a global natural iso}
\label{mk-sec:fbc_gluing}

The affine brick Lemma~\ref{mk-lem:affine_pushforward_pullback_baseChange} produces, over
each affine open \(\operatorname{Spec}(R') \subseteq S'\) (with a chosen affine
\(\operatorname{Spec}(R) \subseteq S\) over its \(g\)-image), an isomorphism
\(e_{R'}\). We glue these into the global natural isomorphism
\(e : g^* \circ f_* \cong f'_* \circ (g')^*\) of the active route. The construction
mirrors the worked sheaf-of-modules descent already carried out in the project's
line-bundle dual, and proceeds in four steps: an affine-restriction naturality lemma
(the one genuinely new statement), a gluing to a sheaf of abelian groups, the
re-imposition of \(\mathcal{O}_{S'}\)-linearity, and the upgrade to an isomorphism
packaged in \(\mathcal{F}\). It rests on the following Mathlib anchors.

\begin{lemma}[Morphism into a sheaf from its restriction to a basis]
  \label{mk-lem:restrictHomEquivHom_mathlib}
  \dcref{custom}

  Let \(F\) be a presheaf and \(F'\) a sheaf (valued in a category) on a space \(X\),
  and let \(\{B_i\}_{i}\) be opens whose range is a basis of the topology of \(X\).
  Then restriction along the basis is a bijection
  \[
    \bigl(B^{\mathrm{op}*}F \longrightarrow B^{\mathrm{op}*}F'\bigr)
    \;\simeq\;
    \bigl(F \longrightarrow F'\bigr),
  \]
  so a morphism into the sheaf \(F'\) is determined by, and constructible from, a
  compatible family of restrictions to the basis opens (with
  \(\operatorname{extend\_hom\_app}\) recovering the component at each \(B_i\)).
\end{lemma}

\begin{lemma}[Re-imposing module linearity on a sheaf morphism]
  \label{mk-lem:presheafOfModules_homMk_mathlib}
  \dcref{custom}

  A morphism \(\varphi\) of the underlying abelian-group presheaves of two presheaves
  of \(R\)-modules \(M_1, M_2\), satisfying the sectionwise linearity equation
  \(\varphi(r \cdot m) = r \cdot \varphi(m)\) for all opens, all scalars
  \(r \in R(U)\) and all sections \(m\), assembles into a morphism
  \(M_1 \to M_2\) of presheaves of \(R\)-modules.
\end{lemma}

\begin{lemma}[Packaging components into a natural isomorphism]
  \label{mk-lem:natIso_ofComponents_mathlib}
  \dcref{custom}

  A family of isomorphisms \(\{\eta_X : G_1(X) \cong G_2(X)\}_X\), natural in \(X\),
  assembles into a natural isomorphism \(G_1 \cong G_2\) of functors.
\end{lemma}
\begin{lemma}[Compatibility of cancellation with localizing the outer base]\label{mk-lem:cancelBaseChange_localization_compat}
  \dcref{custom}

  \uses{mk-lem:cancelBaseChange_mathlib}
  \textit{Source: bespoke project infrastructure; no external citation.}
  Let \(R \to S \to S'\) be a tower of commutative-ring homomorphisms (\(S \to S'\)
  \emph{arbitrary} --- no localization hypothesis), let \(R \to A\) be a further homomorphism, and
  let \(M\) be an \(A\)-module. Realizing the pushouts as the concrete tensor products
  \(A \otimes_R S\) and \(A \otimes_R S'\) (\(A\) the first factor), the cancellation
  isomorphisms \(\operatorname{cancelBaseChange}\) (Lemma~\ref{mk-lem:cancelBaseChange_mathlib}) for the
  \(S\)-tower and the \(S'\)-tower are intertwined by the outer base-change \(S \to S'\): the square
  \[
    \begin{array}{ccc}
      (A \otimes_R S) \otimes_A M & \xrightarrow{\;\operatorname{cancelBaseChange}_S\;} & S \otimes_R M\\
      \big\downarrow{\scriptstyle\,\mathrm{rTensor}\,(\mathrm{id}_A \otimes (S\to S'))} & & \big\downarrow{\scriptstyle\,\mathrm{rTensor}\,(S \to S')}\\
      (A \otimes_R S') \otimes_A M & \xrightarrow{\;\operatorname{cancelBaseChange}_{S'}\;} & S' \otimes_R M
    \end{array}
  \]
  commutes. Equivalently \(\operatorname{cancelBaseChange}\) commutes with changing the outer base
  \(S \to S'\). This is pure \(\operatorname{ModuleCat}\)-level commutative algebra, carrying no
  scheme scaffolding; the localization instance \(S' = S[1/b]\) (with the standard
  \(\mathtt{Algebra}\)/\(\mathtt{IsScalarTower}\) data) is the case consumed by the
  affine-restriction naturality square
  (Lemma~\ref{mk-lem:pushforwardPullbackBaseChange_restrict_naturality}).
\end{lemma}

\begin{proof}
  \uses{mk-lem:cancelBaseChange_mathlib}
  Both \(\operatorname{cancelBaseChange}\) maps are given on simple tensors by the shared formula
  \((a \otimes s) \otimes m \mapsto s \otimes (a \cdot m)\) (Lemma~\ref{mk-lem:cancelBaseChange_mathlib}).
  The square is checked pointwise by a double \(\mathtt{TensorProduct.induction\_on}\) (outer over
  \((A \otimes_R S) \otimes_A M\), inner over \(A \otimes_R S\)); the \(\mathtt{tmul}\) leaf
  reduces, on both routes around the square, to
  \((\operatorname{algebraMap}\,S\,S')(s) \otimes (a \cdot m)\) by the simple-tensor formula together
  with \(\mathrm{rTensor}\) functoriality and the \(A\)-algebra map \(\mathrm{id}_A \otimes (S\to S')\).
  No \(\mathtt{IsLocalizedModule}\) datum or universal property of localized modules is required: the
  identity holds for the \emph{arbitrary} ring map \(S \to S'\). (This is \emph{stronger} than the
  localization-only framing originally planned, and removes
  \(\mathtt{tildeRestriction\_isLocalizedModule}\) from this lemma's cone.)
\end{proof}

\begin{lemma}[Cancellation--base-change compatibility, inverse form]\label{mk-lem:cancelBaseChange_localization_compat_symm}
  \dcref{custom}

  \uses{mk-lem:cancelBaseChange_localization_compat}
  \textit{Source: bespoke project infrastructure; no external citation.}
  In the setting of Lemma~\ref{mk-lem:cancelBaseChange_localization_compat}, the inverse square
  commutes as well:
  \(\mathrm{rTensor}\,(S \to S') \circ \operatorname{cancelBaseChange}_S^{-1}
   = \operatorname{cancelBaseChange}_{S'}^{-1} \circ \mathrm{rTensor}\,(\mathrm{id}_A \otimes (S \to S'))\).
  This is the form consumed by the affine brick, since
  \(\mathtt{baseChangeCancelModuleIso} = \operatorname{cancelBaseChange}^{-1}\).
\end{lemma}
\begin{proof}
  \uses{mk-lem:cancelBaseChange_localization_compat}
  Apply \(\operatorname{cancelBaseChange}_{S'}\) (an isomorphism, hence injective) to both sides and
  rewrite by Lemma~\ref{mk-lem:cancelBaseChange_localization_compat} and
  \(\mathtt{LinearEquiv.apply\_symm\_apply}\) (twice); both sides become equal.
\end{proof}

\subsection{The chart-tower indexing datum (scaffolding)}

The affine-restriction naturality lemma compares the two affine bricks attached to an
inclusion \(\operatorname{Spec}(R'') \subseteq \operatorname{Spec}(R')\) of affine opens
over a common chart \(\operatorname{Spec}(R) \subseteq S\). We package the algebraic datum
of such an inclusion --- four rings \(R, A, R', R''\), the chart map \(\varphi : R \to A\),
the base map \(\psi : R \to R'\), the tower map \(j : R' \to R''\), and the two affine
pushout squares realizing the charts \(B' = A \otimes_R R'\), \(B'' = A \otimes_R R''\) ---
as a single structure, so the naturality statement is phrased against it.

\begin{definition}[Base-change chart tower]\label{mk-def:baseChangeChartTower}
  \dcref{custom}

  \uses{mk-lem:affine_pushforward_pullback_baseChange}
  \textit{Source: bespoke project infrastructure; no external citation.}
  A \emph{base-change chart tower} is the datum of commutative rings \(R, A, R', R''\),
  ring maps \(\varphi : R \to A\), \(\psi : R \to R'\), \(j : R' \to R''\), chart rings
  \(B', B''\) with their pushout legs \(\rho', \sigma', \rho'', \sigma''\), and two
  cocartesian-square witnesses \(h' : \mathrm{IsPushout}\,\varphi\,\psi\,\rho'\,\sigma'\)
  and \(h'' : \mathrm{IsPushout}\,\varphi\,(\psi \circ j)\,\rho''\,\sigma''\). It indexes
  the affine bricks of Lemma~\ref{mk-lem:affine_pushforward_pullback_baseChange} over the
  inclusion \(\operatorname{Spec}(R'') \subseteq \operatorname{Spec}(R')\).
\end{definition}

\begin{definition}[Chart comparison of a tower]\label{mk-def:baseChangeChartTower_connect}
  \dcref{custom}

  \uses{mk-def:baseChangeChartTower}
  \textit{Source: bespoke project infrastructure; no external citation.}
  The canonical ring map \(B' \to B''\) obtained from the universal property of the
  pushout \(B'\) applied to the cocone \((\rho'', j \circ \sigma'')\); it is the affine
  model of the inclusion \(X_{\operatorname{Spec}(R'')} \subseteq X_{\operatorname{Spec}(R')}\).
\end{definition}

\begin{lemma}[Leg compatibilities of the chart comparison]\label{mk-lem:baseChangeChartTower_inl_connect}
  \dcref{custom}

  \uses{mk-def:baseChangeChartTower_connect}
  \textit{Source: bespoke project infrastructure; no external citation.}
  The chart comparison satisfies \(\rho' \circ \mathrm{connect} = \rho''\) (and, by
  Lemma~\ref{mk-lem:baseChangeChartTower_inr_connect}, \(\sigma' \circ \mathrm{connect}
  = j \circ \sigma''\)); both are immediate from the pushout universal property
  (\(\mathrm{inl\_desc}\), \(\mathrm{inr\_desc}\)).
\end{lemma}

\begin{lemma}[Second leg compatibility of the chart comparison]\label{mk-lem:baseChangeChartTower_inr_connect}
  \dcref{custom}

  \uses{mk-def:baseChangeChartTower_connect}
  \textit{Source: bespoke project infrastructure; no external citation.}
  \(\sigma' \circ \mathrm{connect} = j \circ \sigma''\), from the pushout
  universal property.
\end{lemma}

\begin{definition}[The two bricks of a tower]\label{mk-def:baseChangeChartTower_bricks}
  \dcref{custom}

  \uses{mk-def:baseChangeChartTower, mk-lem:affine_pushforward_pullback_baseChange}
  \textit{Source: bespoke project infrastructure; no external citation.}
  The brick \(e_{R'}\) is \(\mathrm{affinePushforwardPullbackBaseChange}\) over
  \(\operatorname{Spec}(R')\) (base map \(\psi\)); the companion brick \(e_{R''}\)
  (\(\mathrm{brickR''}\)) is the same over \(\operatorname{Spec}(R'')\) with the
  \emph{composite} base map \(\psi \circ j\). These are exactly the two isomorphisms the
  restriction-naturality square of
  Lemma~\ref{mk-lem:pushforwardPullbackBaseChange_restrict_naturality} compares.
\end{definition}

\begin{definition}[The second brick of a tower]\label{mk-def:baseChangeChartTower_brickRpp}
  \dcref{custom}

  \uses{mk-def:baseChangeChartTower, mk-lem:affine_pushforward_pullback_baseChange}
  \textit{Source: bespoke project infrastructure; no external citation.}
  The brick \(e_{R''} = \mathrm{affinePushforwardPullbackBaseChange}\) over
  \(\operatorname{Spec}(R'')\) with composite base map \(\psi \circ j\); paired with
  \(\mathrm{brickR'}\) (Definition~\ref{mk-def:baseChangeChartTower_bricks}) in the
  restriction-naturality square.
\end{definition}

\begin{lemma}[Affine-restriction naturality of the cancellation chain]
  \label{mk-lem:pushforwardPullbackBaseChange_restrict_naturality}
  \dcref{custom}

  \uses{mk-lem:affine_pushforward_pullback_baseChange,
        mk-lem:tildeRestriction_isLocalizedModule, mk-lem:cancelBaseChange_mathlib,
        mk-lem:cancelBaseChange_localization_compat, mk-lem:cancelBaseChange_localization_compat_symm,
        mk-def:baseChangeChartTower, mk-def:baseChangeChartTower_bricks}
  Let \(\operatorname{Spec}(R'') \subseteq \operatorname{Spec}(R')\) be an inclusion of
  affine opens of the base \(S'\) (with the same chosen affine
  \(\operatorname{Spec}(R) \subseteq S\) over the \(g\)-image, so that
  \(R \to R' \to R''\) is a tower). Then the affine base-change isomorphisms
  \(e_{R'}\) and \(e_{R''}\) of
  Lemma~\ref{mk-lem:affine_pushforward_pullback_baseChange} are intertwined by the
  structure-sheaf restriction maps: the square with horizontal maps \(e_{R'}, e_{R''}\)
  and vertical maps the restrictions
  \(\operatorname{Spec}(R') \to \operatorname{Spec}(R'')\) on
  \(g^*(f_*\widetilde{M})\) and on \(f'_*((g')^*\widetilde{M})\) commutes. Equivalently,
  the family \(\{e_{R'}\}\) is a morphism of the restricted presheaves over the affine
  basis. This is the one genuinely new lemma of the active route: a concrete
  localization compatibility, replacing the intractable mate coherence.
\end{lemma}

\begin{proof}
  \uses{mk-lem:affine_pushforward_pullback_baseChange,
        mk-lem:tildeRestriction_isLocalizedModule, mk-lem:cancelBaseChange_mathlib,
        mk-lem:cancelBaseChange_localization_compat}
  Each \(e_{R'}\) is the three-move composite of
  Lemma~\ref{mk-lem:affine_pushforward_pullback_baseChange}: two tilde-dictionary
  isomorphisms and \(\widetilde{\operatorname{cancelBaseChange}}\). The restriction
  square factors as the paste of one component square per move, and each commutes for a
  concrete reason. The crucial structural point is that the two bricks live over
  \emph{different base maps}: \(e_{R'}\) is assembled from the dictionaries for the base map
  \(\psi : R \to R'\), whereas \(e_{R''}\) is assembled from the dictionaries for the
  \emph{composite} base map \(\psi \circ j\), with \(j : R' \to R''\) the chart inclusion.
  The vertical restriction maps realize the passage \(\psi \rightsquigarrow \psi\circ j\), so
  the relevant commutation is governed not by fixed-base open-naturality but by the
  \emph{base-map composition cocycle} of the dictionaries.
  \begin{itemize}
    \item \emph{Tilde-dictionary moves.} Restricting from \(\operatorname{Spec}(R')\)
      to \(\operatorname{Spec}(R'')\) replaces each dictionary factor for the base map
      \(\psi\) by the corresponding factor for the composite \(\psi\circ j\). That these
      transitions commute with the dictionary isomorphisms is exactly the base-map
      composition cocycle of the affine dictionaries, instantiated at the tower
      \(R \xrightarrow{\psi} R' \xrightarrow{j} R''\): on the pushforward/\(\Gamma\) side the
      global-sections comparison of Lemma~\ref{mk-lem:gammaPushforwardIso} factors through the
      comparisons for \(\psi\) and \(j\) separately (the source pushforward cocycle absorbed into
      \(\operatorname{Spec}\)-functoriality and the target restriction-of-scalars cocycle folded
      in), and on the pullback side the corresponding factorization of the pullback dictionary of
      Lemma~\ref{mk-lem:pullback_spec_tilde_iso}. (The fixed-base open-naturality of the section
      comparison, recorded in the proof of Lemma~\ref{mk-lem:gammaPushforwardIsoAt}, supplies the
      compatibility of each individual brick with the structure-sheaf restriction maps internal
      to a single base; it is the base-map composition cocycle that reconciles the \emph{two
      distinct base maps} \(\psi\) and \(\psi\circ j\) the square actually compares.)
    \item \emph{Cancellation move.} Localizing the base \(R'\) to \(R''\) carries
      \((R'\otimes_R A)\otimes_A M\) to \((R''\otimes_R A)\otimes_A M = R'' \otimes_{R'}
      \bigl((R'\otimes_R A)\otimes_A M\bigr)\) and \(R'\otimes_R M\) to
      \(R''\otimes_R M = R'' \otimes_{R'} (R'\otimes_R M)\); under these identifications
      \(\operatorname{cancelBaseChange}\) for the tower over \(R''\) is the localization
      of \(\operatorname{cancelBaseChange}\) for the tower over \(R'\). This is exactly the
      cancellation-localization compatibility
      Lemma~\ref{mk-lem:cancelBaseChange_localization_compat} (whose \(R'\)-side localization
      datum is Lemma~\ref{mk-lem:tildeRestriction_isLocalizedModule}).
  \end{itemize}
  Pasting the three component squares yields the restriction square. As the inclusion
  was an arbitrary basic affine inclusion, the family \(\{e_{R'}\}\) is restriction-
  compatible on the affine basis.
\end{proof}

\begin{lemma}[Gluing to a global sheaf-of-abelian-groups morphism]
  \label{mk-lem:pushforwardPullbackBaseChange_abHom}
  \dcref{custom}

  \uses{mk-lem:pushforwardPullbackBaseChange_restrict_naturality,
        mk-lem:restrictHomEquivHom_mathlib, mk-lem:affine_pushforward_pullback_baseChange}
  There is a morphism of sheaves of abelian groups
  \[
    \underline{e} :
    \operatorname{toSheaf}\bigl(g^*(f_*\mathcal{F})\bigr)
    \longrightarrow
    \operatorname{toSheaf}\bigl(f'_*((g')^*\mathcal{F})\bigr)
  \]
  on \(S'\) whose restriction to each affine open
  \(\operatorname{Spec}(R') \subseteq S'\) of the basis is the underlying
  abelian-group morphism of the affine isomorphism \(e_{R'}\).
\end{lemma}

\begin{proof}
  \uses{mk-lem:pushforwardPullbackBaseChange_restrict_naturality,
        mk-lem:restrictHomEquivHom_mathlib, mk-lem:affine_pushforward_pullback_baseChange}
  The affine opens of \(S'\) form a basis of its topology
  (\(X.\operatorname{isBasis\_affineOpens}\)). Forgetting the module structure of both
  \(g^*(f_*\mathcal{F})\) and \(f'_*((g')^*\mathcal{F})\) through
  \(\operatorname{SheafOfModules.toSheaf}\) gives a presheaf source and a \emph{sheaf}
  target of abelian groups. By the restriction naturality
  Lemma~\ref{mk-lem:pushforwardPullbackBaseChange_restrict_naturality}, the affine-local
  components \(\{e_{R'}\}\) (Lemma~\ref{mk-lem:affine_pushforward_pullback_baseChange})
  assemble into a morphism of the basis-restricted presheaves
  \(B^{\mathrm{op}*}F \to B^{\mathrm{op}*}F'\). Mathlib's
  Lemma~\ref{mk-lem:restrictHomEquivHom_mathlib} (basis form) promotes this basis-local
  family to a global morphism of sheaves of abelian groups \(\underline{e}\), with the
  component at each basis open \(\operatorname{Spec}(R')\) recovered as the
  abelian-group morphism of \(e_{R'}\); the corresponding extension property
  \(\operatorname{hom\_ext}\) makes \(\underline{e}\) the unique such global morphism.
\end{proof}

\begin{lemma}[Re-imposing \(\mathcal{O}_{S'}\)-linearity]
  \label{mk-lem:pushforwardPullbackBaseChange_linHom}
  \dcref{custom}

  \uses{mk-lem:pushforwardPullbackBaseChange_abHom,
        mk-lem:presheafOfModules_homMk_mathlib,
        mk-lem:affine_pushforward_pullback_baseChange}
  The abelian-group sheaf morphism \(\underline{e}\) of
  Lemma~\ref{mk-lem:pushforwardPullbackBaseChange_abHom} underlies a morphism of
  \(\mathcal{O}_{S'}\)-modules
  \[
    e^{\mathrm{lin}} :
    g^*(f_*\mathcal{F}) \longrightarrow f'_*((g')^*\mathcal{F}).
  \]
\end{lemma}

\begin{proof}
  \uses{mk-lem:pushforwardPullbackBaseChange_abHom,
        mk-lem:presheafOfModules_homMk_mathlib,
        mk-lem:affine_pushforward_pullback_baseChange}
  The linearity equation \(\underline{e}(r \cdot m) = r \cdot \underline{e}(m)\) is a
  sectionwise identity. It holds on each affine open
  \(\operatorname{Spec}(R') \subseteq S'\) of the basis, because there \(\underline{e}\)
  agrees with the underlying map of \(e_{R'}\), which is
  \(\mathcal{O}_{S'}(\operatorname{Spec}(R'))\)-linear by construction
  (Lemma~\ref{mk-lem:affine_pushforward_pullback_baseChange}). Since the affine opens form
  a basis and both sides of the equation are sections of a sheaf, the identity that
  holds on the basis holds on every open by separatedness. Mathlib's
  Lemma~\ref{mk-lem:presheafOfModules_homMk_mathlib} packages the abelian-group morphism
  \(\underline{e}\) together with this sectionwise linearity witness into the morphism
  \(e^{\mathrm{lin}}\) of \(\mathcal{O}_{S'}\)-modules.
\end{proof}

\begin{lemma}[The global base-change natural isomorphism \(e\)]
  \label{mk-lem:pushforwardPullbackBaseChange_natIso}
  \dcref{custom}

  \uses{mk-lem:pushforwardPullbackBaseChange_linHom,
        mk-lem:modules_isIso_iff_affineOpens,
        mk-lem:affine_pushforward_pullback_baseChange,
        mk-lem:natIso_ofComponents_mathlib}
  The morphism \(e^{\mathrm{lin}}\) of
  Lemma~\ref{mk-lem:pushforwardPullbackBaseChange_linHom} is an isomorphism, and the
  family over quasi-coherent \(\mathcal{F}\) assembles into a natural isomorphism of
  functors
  \[
    e : g^* \circ f_* \;\cong\; f'_* \circ (g')^* .
  \]
  This \(e\) is the active-route deliverable consumed downstream (it commutes with the
  restriction maps by construction); it is built without ever forming the canonical
  adjoint mate.
\end{lemma}

\begin{proof}
  \uses{mk-lem:pushforwardPullbackBaseChange_linHom,
        mk-lem:modules_isIso_iff_affineOpens,
        mk-lem:affine_pushforward_pullback_baseChange,
        mk-lem:natIso_ofComponents_mathlib}
  By the affine-open locality criterion
  Lemma~\ref{mk-lem:modules_isIso_iff_affineOpens}, \(e^{\mathrm{lin}}\) is an isomorphism
  if and only if it restricts to an isomorphism on the sections over every affine open
  of \(S'\). On each such affine open \(\operatorname{Spec}(R')\) it agrees with
  \(e_{R'}\), which is an isomorphism by
  Lemma~\ref{mk-lem:affine_pushforward_pullback_baseChange}; hence \(e^{\mathrm{lin}}\) is
  an isomorphism. Naturality in \(\mathcal{F}\) is checked affine-locally: the brick
  Lemma~\ref{mk-lem:affine_pushforward_pullback_baseChange} is natural in the module
  \(M\), so the components are natural in \(\mathcal{F}\), and this naturality propagates
  through the gluing (the global morphism is the unique extension of natural affine
  data). Finally Mathlib's Lemma~\ref{mk-lem:natIso_ofComponents_mathlib} packages the
  per-\(\mathcal{F}\) isomorphisms into the natural isomorphism \(e\).
\end{proof}

\section{Mathlib stubs imported for the Grassmannian/Quot descent}
\label{mk-sec:fbc_gr_union_stubs}

\begin{lemma}[The pullback--pushforward conjugate sends the pullback coherence to the pushforward coherence]
  \label{mk-lem:conjugateEquiv_pullbackComp_inv_mathlib}
  \dcref{custom}

  Under the conjugation (mate) bijection between the composite pullback--pushforward
  adjunction \((g^{*}\dashv g_{*}) . (f^{*}\dashv f_{*})\) and the adjunction for the
  composite \(\bigl((f\circ g)^{*} \dashv (f\circ g)_{*}\bigr)\), the inverse
  pullback-composition coherence is carried to the pushforward-composition coherence.
\end{lemma}

\begin{lemma}[Sheaf gluing: a compatible family has a unique amalgamation]
  \label{mk-lem:sheaf_existsUnique_gluing_mathlib}
  \dcref{custom}

  Let \(\mathcal{G}\) be a sheaf on \(X\), \(U : \iota \to \mathrm{Opens}(X)\) a family covering an
  open \(V\), and \((s_i)_i\) a compatible family of sections. Then there is a unique
  \(s \in \mathcal{G}(V)\) restricting to each \(s_i\) --- the gluing half of the sheaf condition.
\end{lemma}
